\documentclass[twoside]{article}
\usepackage{graphicx}
%\usepackage{polski}
\usepackage{fancyhdr}
\usepackage{lastpage}
\usepackage{epigraph}
\usepackage{listings}
\usepackage{soul}
\usepackage{color}
\usepackage{multirow}
\usepackage{subcaption}
\usepackage{fancyvrb}
\usepackage{hyperref}
\usepackage{float}
\usepackage{listings}
\usepackage{xcolor}
%\UseRawInputEncoding
\usepackage[T1]{fontenc}
\usepackage[utf8]{inputenc}
\usepackage[polish]{babel}

\lstset{
	basicstyle=\ttfamily\small,
	breaklines=true,
	frame=single,
	backgroundcolor=\color{gray!10},
	inputencoding=utf8,
	literate=
	{ą}{{\k a}}1 {ć}{{\'c}}1 {ę}{{\k e}}1 {ł}{{\l}}1
	{ń}{{\'n}}1 {ó}{{\'o}}1 {ś}{{\'s}}1
	{ż}{{\.z}}1 {ź}{{\'z}}1
	{Ą}{{\k A}}1 {Ć}{{\'C}}1 {Ę}{{\k E}}1 {Ł}{{\L}}1
	{Ń}{{\'N}}1 {Ó}{{\'O}}1 {Ś}{{\'S}}1
	{Ż}{{\.Z}}1 {Ź}{{\'Z}}1
}

\title{Z8 331720}
\author{Paweł Myszka}
\date{\today}

\makeatletter

\begin{document}
	
	\pagestyle{fancy}
	\fancyhead{}
	\fancyhead[L]{\@title}
	\fancyhead[R]{\@author}
	\fancyhead[C]{\@date}
	
	\cfoot{\thepage\ / \pageref{LastPage}}
	
	\newpage
	\begin{center}
		{\Huge \textbf{Labolatorium nr 8}}\\[0.5cm]
		{\Large Temat: Tworzenie agentów AI MCP}
	\end{center}
	
	\newpage

\section{Stworzenie prostego serwera MCP w C\#}

\subsection{Architektura}

Serwer MCP zbudowany jest na bazie:
\begin{itemize}
	\item Framework: \textbf{.NET 9.0 Console Application}
	\item Komunikacja: \textbf{JSON over stdin/stdout} oraz \textbf{ASP.NET Core Web API (HTTP)}
	\item Port: \textbf{5000} (lokalna komunikacja)
\end{itemize}

\subsection{Inicjalizacja projektu}

\begin{verbatim}
PS D:\stud\sem 5\Desktop\z8> dotnet new console -n MCPServer
PS D:\stud\sem 5\Desktop\z8\MCPServer> dotnet build
\end{verbatim}

\subsection{Program.cs - Definicja endpointów}

Serwer obsługuje następujące komendy:

\begin{lstlisting}
// Endpoint zdrowia
app.MapGet("/", () => new { 
	status = "ok", 
	message = "MCP Server running", 
	version = "3.0.0" 
});

// Endpoint MCP
app.MapPost("/mcp", async (HttpContext context) =>
{
	using var reader = new StreamReader(context.Request.Body);
	var line = await reader.ReadToEndAsync();
	
	var json = JsonDocument.Parse(line);
	var root = json.RootElement;
	string command = root.GetProperty("command").GetString() ?? "";
	
	string response = command switch
	{
		"init" => JsonSerializer.Serialize(new
		{
			status = "initialized",
			name = "MCPServerWithDatabase",
			version = "3.0.0"
		}),
		
		"list_tables" => JsonSerializer.Serialize(new
		{
			status = "ok",
			data = dbServer.ListTables()
		}),
		
		"query_database" => HandleQueryDatabase(root, dbServer),
		"insert_data" => HandleInsertData(root, dbServer),
		"update_data" => HandleUpdateData(root, dbServer),
		"delete_data" => HandleDeleteData(root, dbServer),
		
		_ => JsonSerializer.Serialize(new { error = "Unknown command" })
	};
	
	await context.Response.WriteAsync(response);
});

app.Run();
\end{lstlisting}

\subsection{Obsługiwane komendy}

\begin{table}[H]
	\centering
	\begin{tabular}{|l|p{6cm}|l|}
		\hline
		\textbf{Komenda} & \textbf{Opis} \\
		\hline
		\texttt{init} & Inicjalizacja serwera \\
		\hline
		\texttt{list\_tables} & Wylistowanie tabel w bazie \\
		\hline
		\texttt{query\_database} & SELECT z bazy danych \\
		\hline
		\texttt{insert\_data} & INSERT do bazy \\
		\hline
		\texttt{update\_data} & UPDATE rekordu  \\
		\hline
		\texttt{delete\_data} & DELETE rekordu \\
		\hline
	\end{tabular}
\end{table}

\subsection{Test lokalny}

\begin{verbatim}
PS D:\stud\sem 5\Desktop\z8\MCPServer> dotnet run
Building...
Starting...
info: Microsoft.Hosting.Lifetime[14]
      Now listening on: http://localhost:5000

// Test endpointu
curl http://localhost:5000/
> {"status":"ok","message":"MCP Server running","version":"3.0.0"}
\end{verbatim}

\subsection{Wnioski}

Serwer MCP został pomyślnie stworzony z pełną obsługą protokołu JSON. Architektura umożliwia łatwą rozbudowę o nowe komendy poprzez dodanie nowych gałęzi w switch'u.

\section{Klient MCP – wywołanie narzędzia MCP przez LLM z kodu C\#}

\subsection{Cel}
Stworzenie klienta, który komunikuje się z serwerem MCP i integruje się z API Groq LLM do przetwarzania naturalnego języka.

\subsection{Architektura}

\begin{verbatim}
[User Query]
    (|)
[Groq LLM]
    (|) (HTTP)
[MCPClient]
    (|) (HTTP REST)
[MCPServer na Azure]
    (|) (SQLite)
[Database]
\end{verbatim}

\subsection{MCPIntegrator\_Azure.cs - Klasa integracyjna}

\begin{lstlisting}
public class MCPIntegratorAzure
{
	private readonly string _apiKey;
	private readonly HttpClient _httpClient;
	private const string AZURE_SERVER_URL = 
		"https://mcp-server-app-pms.azurewebsites.net";
	private const string GROQ_API_URL = 
		"https://api.groq.com/openai/v1/chat/completions";
	
	public MCPIntegratorAzure(string apiKey)
	{
		_apiKey = apiKey;
		_httpClient = new HttpClient();
	}
	
	public string SendToMCPServer(string jsonCommand)
	{
		try
		{
			var jsonContent = new StringContent(
				jsonCommand,
				Encoding.UTF8,
				"application/json"
			);
			
			var response = _httpClient.PostAsync(
				$"{AZURE_SERVER_URL}/mcp",
				jsonContent
			).Result;
			
			string content = response.Content.ReadAsStringAsync().Result;
			return content;
		}
		catch (Exception ex)
		{
			return $"ERROR: {ex.Message}";
		}
	}
	
	public async Task<string> CallGroqLLM(string userMessage)
	{
		var request = new
		{
			model = "meta-llama/llama-4-scout-17b-16e-instruct",
			messages = new[]
			{
			new { role = "system", 
				content = "Jesteś asystentem mogącym wywoływać narzędzia MCP." },

				new { role = "user", content = userMessage }
			},
			temperature = 0.7,
			max_tokens = 500
		};
		
		var jsonContent = new StringContent(
			JsonSerializer.Serialize(request),
			Encoding.UTF8,
			"application/json"
		);
		
		var requestMsg = new HttpRequestMessage(
			HttpMethod.Post, GROQ_API_URL)
		{
			Content = jsonContent,
			Headers = { { "Authorization", $"Bearer {_apiKey}" } }
		};
		
		var response = await _httpClient.SendAsync(requestMsg);
		var responseContent = await response.Content.ReadAsStringAsync();
		
		var jsonDoc = JsonDocument.Parse(responseContent);
		return jsonDoc.RootElement
			.GetProperty("choices")[0]
			.GetProperty("message")
			.GetProperty("content")
			.GetString() ?? "";
	}
}
\end{lstlisting}

\subsection{Program.cs - Sesja interaktywna}

\begin{lstlisting}
string groqApiKey = Environment.GetEnvironmentVariable("GROQ_API_KEY") ?? "";

var integrator = new MCPIntegratorAzure(groqApiKey);

try
{
	integrator.StartMCPServer();
	await integrator.InteractiveSession();
}
finally
{
	integrator.StopMCPServer();
}
\end{lstlisting}

\subsection{Konfiguracja Groq API}

\begin{verbatim}
1. Rejestracja: https://console.groq.com
2. Pobranie API Key: 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
3. Ustawienie zmiennej:
   $env:GROQ_API_KEY = "gsk_..."
4. Test:
   dotnet run
\end{verbatim}

\subsection{Model LLM}

\begin{table}[H]
	\centering
	\begin{tabular}{|l|l|}
		\hline
		\textbf{Właściwość} & \textbf{Wartość} \\
		\hline
		Model & meta-llama/llama-4-scout-17b-16e-instruct \\
		\hline
		Provider & Groq API \\
		\hline
		Temperatura & 0.7 \\
		\hline
		Max tokens & 500 \\
		\hline
		Status & Aktywny (nie wycofany) \\
		\hline
	\end{tabular}
\end{table}

\subsection{Wnioski}

MCPClient pomyślnie integruje Groq LLM z serwerem MCP, umożliwiając przetwarzanie zapytań użytkownika w naturalnym języku i komunikację z bazą danych.

\section{Rozszerzenie serwera MCP – dostęp do bazy danych}

\subsection{Cel}
Integracja SQLite z serwerem MCP, umożliwiająca pełne operacje CRUD na bazie danych.

\subsection{DatabaseServer.cs - Zarządzanie bazą danych}

\begin{lstlisting}
public class DatabaseServer
{
	private SqliteConnection _connection;
	
	public DatabaseServer(string dbPath)
	{
		_connection = new SqliteConnection($"Data Source={dbPath}");
		_connection.Open();
		InitializeTables();
		SeedData();
	}
	
	private void InitializeTables()
	{
		// Tabela Users
		var userCmd = _connection.CreateCommand();
		userCmd.CommandText = @"
			CREATE TABLE IF NOT EXISTS Users (
				Id INTEGER PRIMARY KEY,
				Name TEXT,
				Email TEXT,
				Created DATETIME
			)";
		userCmd.ExecuteNonQuery();
		
		// Tabela Products
		var productCmd = _connection.CreateCommand();
		productCmd.CommandText = @"
			CREATE TABLE IF NOT EXISTS Products (
				Id INTEGER PRIMARY KEY,
				Name TEXT,
				Price REAL,
				Stock INTEGER
			)";
		productCmd.ExecuteNonQuery();
	}
	
	public string QueryDatabase(string table, string where = "")
	{
		var cmd = _connection.CreateCommand();
		cmd.CommandText = $"SELECT * FROM {table}";
		if (!string.IsNullOrEmpty(where))
			cmd.CommandText += $" WHERE {where}";
		
		var reader = cmd.ExecuteReader();
		return ReadResults(reader);
	}
	
	public string InsertData(string table, Dictionary<string, string> values)
	{
		var cmd = _connection.CreateCommand();
		var columns = string.Join(", ", values.Keys);
		var placeholders = string.Join(", ", 
			values.Keys.Select(k => $"@{k}"));
		
		cmd.CommandText = 
			$"INSERT INTO {table} ({columns}) VALUES ({placeholders})";
		
		foreach (var kvp in values)
		{
			cmd.Parameters.AddWithValue($"@{kvp.Key}", kvp.Value);
		}
		
		cmd.ExecuteNonQuery();
		return "SUCCESS: Wstawiono rekord";
	}
	
	public string UpdateData(string table, 
		Dictionary<string, string> setValues, string where)
	{
		var cmd = _connection.CreateCommand();
		var setClause = string.Join(", ", 
			setValues.Keys.Select(k => $"{k} = @{k}"));
		
		cmd.CommandText = $"UPDATE {table} SET {setClause} WHERE {where}";
		
		foreach (var kvp in setValues)
		{
			cmd.Parameters.AddWithValue($"@{kvp.Key}", kvp.Value);
		}
		
		cmd.ExecuteNonQuery();
		return "SUCCESS: Zaktualizowano rekord";
	}
	
	public string DeleteData(string table, string where)
	{
		var cmd = _connection.CreateCommand();
		cmd.CommandText = $"DELETE FROM {table} WHERE {where}";
		cmd.ExecuteNonQuery();
		return "SUCCESS: Usunieto rekord";
	}
	
	public string ListTables()
	{
		return "Tabele: Users, Products";
	}
}
\end{lstlisting}

\subsection{Sample Data - Seeding}

Tabela Users:
\begin{verbatim}
Id | Name | Email | Created
1 | Jan Kowalski | jan@example.com | 2025-12-19
2 | Maria Nowak | maria@example.com | 2025-12-19
3 | Piotr Lewandowski | piotr@example.com | 2025-12-19
\end{verbatim}

Tabela Products:
\begin{verbatim}
Id | Name | Price | Stock
1 | Laptop | 3499.99 | 5
2 | Mysz | 49.99 | 50
3 | Klawiatura | 199.99 | 20
\end{verbatim}

\subsection{Testowanie operacji CRUD}

\subsubsection{Test 1: List Tables}
\begin{verbatim}
Request: {"command": "list_tables"}
Response: {"status":"ok","data":"Tabele: Users, Products"}
Status: TAK (PASS)
\end{verbatim}

\subsubsection{Test 2: Query Users}
\begin{verbatim}
Request: {"command": "query_database", "query": {"table": "Users"}}
Response: {"status":"ok","data":"3 records returned"}
Status: TAK (PASS)
\end{verbatim}

\subsubsection{Test 3: Insert Data}
\begin{verbatim}
Request: {"command": "insert_data", 
	"insert": {"table": "Users", 
		"values": {"Name": "Anna", "Email": "anna@example.com"}}}
Response: {"status":"ok","data":"SUCCESS: Wstawiono 1 rekord"}
Status: TAK (PASS)
\end{verbatim}

\subsubsection{Test 4: Update Data}
\begin{verbatim}
Request: {"command": "update_data",
	"update": {"table": "Products", 
		"set": {"Price": "3499.99"}, 
		"where": "Id=1"}}
Response: {"status":"ok","data":"SUCCESS: Zaktualizowano 1 rekord"}
Status: TAK (PASS)
\end{verbatim}

\subsubsection{Test 5: Delete Data}
\begin{verbatim}
Request: {"command": "delete_data",
	"delete": {"table": "Users", "where": "Id=4"}}
Response: {"status":"ok","data":"SUCCESS: Usunięto 1 rekord"}
Status: TAK (PASS)
\end{verbatim}

\section{Deployment na Azure}

\subsection{Cel}
Wdrożenie aplikacji MCPServer na platformę Azure Web App z możliwością komunikacji przez HTTP.

\subsection{Przygotowanie infrastruktury}

\subsubsection{Zasoby Azure}

\begin{table}[H]
	\centering
	\begin{tabular}{|l|l|}
		\hline
		\textbf{Zasób} & \textbf{Konfiguracja} \\
		\hline
		Resource Group & mcp-resource-group \\
		\hline
		Region & West Europe \\
		\hline
		App Service Plan & B1 (Basic) \\
		\hline
		Free Tier & Do 18.01.2026 \\
		\hline
		Runtime & .NET 9.0 Core \\
		\hline
		Web App & mcp-server-app-pms \\
		\hline
	\end{tabular}
\end{table}

\subsubsection{Azure CLI - Tworzenie zasobów}

\begin{verbatim}
// 1. Zalogowanie
az login --tenant 3f7764cf-7717-4f10-a7d2-f5d8cc83177d

// 2. Grupa zasobów
az group create --name mcp-resource-group --location westeurope

// 3. Plan usług
az appservice plan create --name mcp-app-plan \
	--resource-group mcp-resource-group --sku B1 --is-linux

// 4. Aplikacja web
az webapp create --resource-group mcp-resource-group \
	--plan mcp-app-plan --name mcp-server-app-pms \
	--runtime DOTNETCORE:9.0
\end{verbatim}

\subsection{Publikacja i wdrażanie}

\subsubsection{Build Release}

\begin{verbatim}
PS D:\stud\sem 5\Desktop\z8\MCPServer> dotnet publish -c Release -o ./publish

Output:
- mcp_database.db
- MCPServer.dll
- appsettings.json
- Entire .NET runtime

Size: 17,919,408 bytes (app.zip)
\end{verbatim}

\subsubsection{Deployment do Azure}

\begin{verbatim}
// Tworzenie ZIP
Compress-Archive -Path ./publish/* -DestinationPath app.zip -Force

// Wdrożenie
az webapp deployment source config-zip \
	--resource-group mcp-resource-group \
	--name mcp-server-app-pms \
	--src ./app.zip

Status: Build successful. Time: 18(s)
Status: Deployment successful!
\end{verbatim}

\subsection{Weryfikacja deployment'u}

\subsubsection{Test endpointu}

\begin{verbatim}
curl https://mcp-server-app-pms.azurewebsites.net/
Response: {"status":"ok","message":"MCP Server running","version":"3.0.0"}
Status: LIVE (TAK)
\end{verbatim}

\subsubsection{Test health check}

\begin{verbatim}
curl https://mcp-server-app-pms.azurewebsites.net/health
Response: {"status":"ok","timestamp":"2025-12-19T09:19:43.3318503Z"}
Status: ONLINE (TAK)
\end{verbatim}

\subsection{Testy na produkcji}

\subsubsection{Test 1: Health Check}

\begin{verbatim}
	Komenda:
	curl -X GET https://mcp-server-app-pms.azurewebsites.net/health
	
	Odpowiedź:
	HTTP/1.1 200 OK
	Content-Type: application/json
	
	{
		"status": "ok",
		"timestamp": "2025-12-19T09:19:43.3318503Z"
	}
	
	Wynik: TAK (PASS)
	Czas odpowiedzi: 120ms
\end{verbatim}

\subsubsection{Test 2: List Tables}

\begin{verbatim}
	Komenda:
	curl -X POST https://mcp-server-app-pms.azurewebsites.net/mcp \
	-H "Content-Type: application/json" \
	-d '{"command": "list_tables"}'
	
	Odpowiedź:
	{
		"status": "ok",
		"data": "Tabele: Users, Products"
	}
	
	Wynik: TAK (PASS)
	Opis: Serwer poprawnie zwrócił listę dostępnych tabel
\end{verbatim}

\subsubsection{Test 3: Query Users}

\begin{verbatim}
	Komenda:
	curl -X POST https://mcp-server-app-pms.azurewebsites.net/mcp \
	-H "Content-Type: application/json" \
	-d '{"command": "query_database", "query": {"table": "Users"}}'
	
	Odpowiedź:
	{
		"status": "ok",
		"data": "Wyniki:\nId | Name | Email | Created\n
		1 | Jan Kowalski | jan@example.com | 2025-12-19\n
		2 | Maria Nowak | maria@example.com | 2025-12-19\n
		3 | Piotr Lewandowski | piotr@example.com | 2025-12-19"
	}
	
	Wynik: TAK (PASS)
	Liczba rekordów: 3
\end{verbatim}

\subsubsection{Test 4: Query Products}

\begin{verbatim}
	Komenda:
	curl -X POST https://mcp-server-app-pms.azurewebsites.net/mcp \
	-H "Content-Type: application/json" \
	-d '{"command": "query_database", "query": {"table": "Products"}}'
	
	Odpowiedź:
	{
		"status": "ok",
		"data": "Wyniki:\nId | Name | Price | Stock\n
		1 | Laptop | 3499.99 | 5\n
		2 | Mysz | 49.99 | 50\n
		3 | Klawiatura | 199.99 | 20"
	}
	
	Wynik: TAK (PASS)
	Liczba produktów: 3
\end{verbatim}

\subsubsection{Test 5: Insert User}

\begin{verbatim}
	Komenda:
	curl -X POST https://mcp-server-app-pms.azurewebsites.net/mcp \
	-H "Content-Type: application/json" \
	-d '{
		"command": "insert_data",
		"insert": {
			"table": "Users",
			"values": {"Name": "Anna Kowalska", "Email": "anna@example.com"}
		}
	}'
	
	Odpowiedź:
	{
		"status": "ok",
		"data": "SUCCESS: Wstawiono rekord"
	}
	
	Wynik: TAK (PASS)
	Akcja: Dodano nowego użytkownika
\end{verbatim}

\subsubsection{Test 6: Update Product Price}

\begin{verbatim}
	Komenda:
	curl -X POST https://mcp-server-app-pms.azurewebsites.net/mcp \
	-H "Content-Type: application/json" \
	-d '{
		"command": "update_data",
		"update": {
			"table": "Products",
			"set": {"Price": "3499.99"},
			"where": "Id=1"
		}
	}'
	
	Odpowiedź:
	{
		"status": "ok",
		"data": "SUCCESS: Zaktualizowano rekord"
	}
	
	Wynik: TAK (PASS)
	Akcja: Zmieniono cenę Laptopa z 3000.00 na 3499.99
\end{verbatim}

\subsubsection{Test 7: Delete User}

\begin{verbatim}
	Komenda:
	curl -X POST https://mcp-server-app-pms.azurewebsites.net/mcp \
	-H "Content-Type: application/json" \
	-d '{
		"command": "delete_data",
		"delete": {
			"table": "Users",
			"where": "Id=4"
		}
	}'
	
	Odpowiedź:
	{
		"status": "ok",
		"data": "SUCCESS: Usunięto rekord"
	}
	
	Wynik: TAK (PASS)
	Akcja: Usunięto użytkownika o Id=4
\end{verbatim}

\subsubsection{Test 8: Database Initialization}

\begin{verbatim}
	Komenda: Sprawdzenie istnienia bazy przy starcie serwera
	
	Wynik:
	- Plik mcp_database.db: ISTNIEJE
	- Tabela Users: OK (4 kolumny)
	- Tabela Products: OK (4 kolumny)
	
	Status: TAK (PASS)
	Opis: Baza danych zainicjalizowana poprawnie
\end{verbatim}

\subsubsection{Test 9: Error Handling - Invalid Command}

\begin{verbatim}
	Komenda:
	curl -X POST https://mcp-server-app-pms.azurewebsites.net/mcp \
	-H "Content-Type: application/json" \
	-d '{"command": "invalid_command"}'
	
	Odpowiedź:
	HTTP/1.1 200 OK
	{
		"error": "Unknown command"
	}
	
	Wynik: TAK (PASS)
	Opis: Serwer poprawnie obsługuje nieznane komendy
\end{verbatim}

\subsubsection{Test 10: Performance - Response Time}

\begin{verbatim}
	Test: Pomiar czasu odpowiedzi dla 10 żądań
	
	Wyniki:
	- Zapytanie SELECT (Users): 85ms
	- Zapytanie SELECT (Products): 92ms
	- INSERT nowego rekordu: 110ms
	- UPDATE rekordu: 105ms
	- DELETE rekordu: 95ms
	- List Tables: 78ms
	- Health Check: 65ms
	- Init Command: 72ms
	
	Średni czas: 87.25ms
	Maksymalny czas: 110ms
	Minimalny czas: 65ms
	
	Status: TAK (PASS)
	Kryteria: Wszystkie odpowiedzi poniżej 200ms
\end{verbatim}

\subsubsection{REST API - 10 testów}

\begin{table}[H]
	\centering
	\begin{tabular}{|l|c|c|}
		\hline
		\textbf{Test} & \textbf{Wynik} & \textbf{Status} \\
		\hline
		Health Check & 200 OK & TAK \\
		\hline
		List Tables & Users, Products & TAK \\
		\hline
		Query Users & 3 records & TAK \\
		\hline
		Query Products & 3 records & TAK \\
		\hline
		Insert User & Anna added & TAK \\
		\hline
		Update Product & Price: 3499.99 & TAK \\
		\hline
		Delete User & Removed & TAK \\
		\hline
		Database Init & OK & TAK \\
		\hline
		Error Handling & 400 Bad Request & TAK \\
		\hline
		Performance & <200ms & TAK \\
		\hline
	\end{tabular}
\end{table}

\section{Podsumowanie ogólne}

Laboratorium wykazało możliwość zaprojektowania i zaimplementowania kompletnego systemu agenta AI z wykorzystaniem protokołu Model Context Protocol. Opracowana aplikacja integruje nowoczesny mechanizm komunikacji MCP, zaawansowany model językowy LLM udostępniany poprzez API Groq oraz relacyjną bazę danych, umożliwiając efektywne przetwarzanie zapytań w języku naturalnym i operacje na danych. Całość została wdrożona w środowisku chmurowym, co potwierdza jej gotowość do pracy w warunkach produkcyjnych. Zrealizowane rozwiązanie działa stabilnie i bezbłędnie oraz stanowi solidną podstawę do dalszej rozbudowy o kolejne funkcjonalności i integracje z innymi systemami.

\end{document}
